\subsection{Recoverability and Bayesian Diagnostics}
\label{sec:recoverability_diagnostics}

\paragraph{Benchmarks.}
We evaluate whether BHPOP can recover latent workflow structure from execution traces under
controlled noise and trace diversity.
We report results on (i) a public workflow-DAG benchmark (WFCommons WfFormat instances) and
(ii) synthetic hierarchical partial orders sampled from our prior, which enables sweeps over
trace budget, latent dimension $K$, inverse temperature $\beta$, and trembling-hand noise $\epsilon$.
To avoid leakage from dataset-specific identifiers, we randomize node IDs within each workflow
instance prior to inference (an ablation is provided in Appendix~\ref{app:recoverability_extra}).

\paragraph{Observation model.}
Given a candidate partial order $h$, we score a trace $y_{1:T}$ with the frontier-softmax
queue-jump likelihood from Section~\ref{sec:model}, parameterized by inverse temperature
$\beta$ and trembling-hand mass $\epsilon$.
Intuitively, $\beta$ controls how strongly the agent follows a canonical plan, while $\epsilon$
allocates small probability to off-frontier (infeasible) actions.

\paragraph{Posterior prediction target.}
From post burn-in samples, we estimate posterior cover-edge probabilities
$\hat p_{ij} \triangleq P(i \prec j \mid \mathcal{D})$ on the transitive reduction
(Hasse diagram) and form a point estimate $\widehat{G}$ by thresholding $\hat p_{ij} \ge 0.5$.
Unless otherwise stated, all structural metrics are computed on cover edges to avoid
double-counting implied transitive relations.

\paragraph{Structural recovery.}
Let $\mathrm{TR}(\cdot)$ denote transitive reduction and let $E^\star$ and $\widehat{E}$ be
the directed edge sets of $\mathrm{TR}(G^\star)$ and $\mathrm{TR}(\widehat{G})$.
We report precision/recall/F1 on cover edges and the structural Hamming distance (SHD)
on cover edges:
\begin{equation}
\mathrm{SHD}(\widehat{E},E^\star)
= |\widehat{E}\setminus E^\star| + |E^\star\setminus \widehat{E}|.
\end{equation}
(A reversed edge contributes two operations: one deletion and one addition.)

\paragraph{Predictive fit (WAIC).}
We compute WAIC using the standard pointwise log-likelihood matrix
$\ell_{i,s}=\log p(y_i\mid\theta^{(s)})$ over posterior draws $s$.
To enable fair comparison across different trace budgets, we report a normalized score
\begin{equation}
\mathrm{WAIC/step} \triangleq \frac{\mathrm{WAIC}}{\sum_i |y_i|},
\end{equation}
since unnormalized WAIC scales approximately linearly with the number of observed steps.

\paragraph{Uncertainty calibration.}
We assess whether posterior edge probabilities are meaningful by pooling predicted
probabilities $\hat p_{ij}$ across workflow instances and comparing them to ground-truth edge
indicators.
We report a reliability diagram and expected calibration error (ECE) with $B$ bins:
\begin{equation}
\mathrm{ECE}=\sum_{b=1}^B \frac{|B_b|}{N}\left|\mathrm{acc}(B_b)-\mathrm{conf}(B_b)\right|.
\end{equation}

\paragraph{Main findings.}
In the latest synthetic sweep (Appendix~\ref{app:recoverability_extra}, Table~\ref{tab:synthetic_latest_run}),
structural recovery and uncertainty diagnostics depend on both $\beta$ and the latent poset
dimension $K$.
Because dominance-based partial orders become sparser as $K$ increases, perfect recovery at large
$K$ can correspond to recovering very few edges; we therefore report $|E^\star|$ alongside F1/SHD.

\begin{figure}[t]
\centering
% Replace paths with your paper's figure directory.
\includegraphics[width=0.49\linewidth]{figures/recoverability/f1_vs_traces}
\includegraphics[width=0.49\linewidth]{figures/recoverability/shd_vs_traces}
\caption{\textbf{Structural recoverability.} F1 (left) and SHD (right) on cover edges as a function
of the number of observed traces. Curves show mean $\pm$ std over random seeds.}
\label{fig:recoverability_scaling}
\end{figure}

\begin{figure}[t]
\centering
\includegraphics[width=0.49\linewidth]{figures/recoverability/waic_per_step_vs_traces}
\includegraphics[width=0.49\linewidth]{figures/recoverability/calibration}
\caption{\textbf{Bayesian diagnostics.} Left: predictive fit via WAIC/step (lower is better).
Right: reliability diagram for posterior edge probabilities with ECE reported in the legend.}
\label{fig:recoverability_diagnostics}
\end{figure}
